Este trabajo presenta un análisis empírico y comparativo del rendimiento computacional de algoritmos de ordenamiento (\textit{Merge Sort}, \textit{Quick Sort}, \textit{Patience Sort} y \textit{Sort} default) y multiplicación matricial (algoritmo clásico \textit{Naive} O $\matchcal(N^3)$ frente al método de Strassen O $\mathcal(N^{2.807})$). A través de un marco de pruebas automatizado en C++17 y un pipeline analítico en Python, se evaluaron tiempos de ejecución sobre diversos ordenes de magnitud y orden de distribución. Los resultados confirman predicciones asintóticas teóricas, evidenciando cómo otros factores (caché, asignación de memoria y adaptar la recursión) determinan la viabilidad práctica de los algoritmos yendo más allá de su complejidad teórica matemática.